\section{Method}\label{sec:method}

\subsection{Problem setting}\label{sec:problem}
At frame $t$ a frozen detector emits candidates $D_t=\{(\mathbf{x}_i,s_i,c_i)\}$
with boxes $\mathbf{x}_i$, scores $s_i\in(0,1)$, and class labels $c_i$, down
to an emission floor. The layer works on the logit of each score,
\begin{equation}
z_i=\log\frac{s_i}{1-s_i}.
\label{eq:logit}
\end{equation}
A frozen tracker, called the host, declares a host contract
$h=(a,b,\ell,m_0)$. Here $a$ is the score threshold of its first association,
$b$ is the lowest score from which a new track can start, $\ell$ is the lowest
score that any of its association stages uses, and $m_0$ is its matching
tolerance (one minus its minimum overlap). A single-stage tracker has
$\ell=\min(a,b)$, and a two-stage tracker has $\ell<\min(a,b)$.

For every frame the layer outputs the set $K_t$ of candidates passed to the
tracker, the scores $\tilde{s}$ passed with them, and the thresholds
$(a_t,b_t)$ and tolerance $m_t$ that the tracker uses in that frame:
\begin{equation}
\left(K_t,\tilde{s},a_t,b_t,m_t\right)=
\mathcal{L}\left(Z_{<t},D_t,\mathcal{T}_{t-1},g_t;\,h\right).
\label{eq:layer}
\end{equation}
$Z_{<t}$ are the logits of earlier frames, $\mathcal{T}_{t-1}$ are the
tracks that the tracker reported for frame $t-1$, and $g_t$ is an optional
image-motion cue of frame $t$. The layer works online. Its statistics use only
frames before $t$, and the data of frame $t$ enter its state after the
decision. No future frame is used.

The layer has four requirements: it needs no training and no labelled
calibration data; one frozen configuration serves every tracker, detector, and
dataset; the output stays unchanged when the tracker's operating point already
fits the score stream; and its cost per frame is small next to the detector.

\subsection{Overview}
Figure~\ref{fig:arch} shows \scl{}. It has four parts: stream statistics, a
regime estimate, a candidate policy for each regime, and a continuation
rescue. The host contract enters the rescue and the thresholds passed to the
tracker. The detector and the tracker are not modified.

\begin{figure*}[t]
\centering
\begin{tikzpicture}[
  box/.style={draw, rounded corners=2pt, align=center, font=\footnotesize, minimum height=9mm, minimum width=22mm, inner sep=3pt},
  arr/.style={-{Stealth[length=2mm]}, thick}, lab/.style={font=\scriptsize}]
\node[box] (det) at (0,0) {Frozen\\detector};
\node[box] (pol) at (4.3,0) {Candidate policy\\(clean or noisy)};
\node[box] (res) at (8.5,0) {Continuation\\rescue};
\node[box] (trk) at (12.7,0) {Frozen\\tracker};
\node[box] (stat) at (0,-2.1) {Stream statistics:\\nested Otsu on logits\\of frames $t-W..t-1$};
\node[box] (reg) at (4.3,-2.1) {Regime estimate\\clean if $\bar\rho_t \ge \tfrac12$};
\node[box] (hc) at (10.6,-2.1) {Host contract\\$(a, b, \ell, m_0)$};
\draw[arr] (det) -- node[lab, above]{candidates} (pol);
\draw[arr] (pol) -- (res);
\draw[arr] (res) -- node[lab, above]{$K_t$, scores} (trk);
\draw[arr] (det) -- (stat);
\draw[arr] (stat) -- node[lab, above]{$t_1, t_2, \rho$} (reg);
\draw[arr] (reg) -- (pol);
\draw[arr] (hc.north) ++(-0.6,0) |- ++(0,0.45) -| (res.south);
\draw[arr] (hc.north) ++(0.6,0) |- ++(0,0.45) -| node[lab, pos=0.75, right]{$a_t, b_t, m_t$} (trk.south);
\draw[arr] (trk.north) -- ++(0,0.45) -| node[lab, pos=0.25, above]{tracks of frame $t-1$} (res.north);
\end{tikzpicture}
\caption{Self-calibrating control layer (\scl{}) between a frozen detector and a
frozen tracker. The statistics use earlier frames only. The host contract is the
tracker's published operating point. The tracker receives the selected
candidates, their scores, and its thresholds for the frame.}
\label{fig:arch}
\end{figure*}

\subsection{Self-calibration from the score stream}\label{sec:calib}
The layer pools the logits of the candidates of the previous $W=10$ frames and
splits them twice with the exact Otsu threshold \cite{otsu1979}. The first
split $t_1$ separates background from foreground. A second split of the
foreground, $t_2$, separates ambiguous from confident candidates. The
confident share of the foreground,
\begin{equation}
\rho_t=\frac{\left|\{i: z_i\ge t_2\}\right|}{\left|\{i: z_i\ge t_1\}\right|},
\label{eq:rho}
\end{equation}
summarizes how cleanly the detector separates objects from clutter. Because
the splits are computed on logits, they follow any increasing affine change
of the logit, such as a change of temperature, and the bands keep the same
candidates.

The bands have their intended meaning only if the stream has a background
mode. When the class below $t_1$ holds fewer pooled candidates than the
foreground, which happens when a detector emits few low-scoring candidates, the
first split falls inside the objects. In that case the frame counts as clean
evidence, $\rho_t=1$. This background check was added after a development
case in which sequence starts with about ten candidates per frame produced
splits at $t_1\approx0.6$ and $t_2\approx0.93$ in probability terms and a
wrong noisy decision.

\subsection{Intervention regimes}\label{sec:regimes}
\emph{Regime estimate.} The stream is clean at frame $t$ if the median
$\bar\rho_t$ of $\rho$ over all non-cold frames so far is at least one half,
and noisy otherwise. Before the first valid split, for example at the first
frame of a stream, the regime is cold. The regime is therefore a property of
the detector and scene stream, and a short dip in confidence, for example under fast camera
motion, does not flip it.

\emph{Cold and clean regimes.} The tracker's own thresholds are kept. In the
clean regime they are lowered only as far as $t_2$ when the tracker would
otherwise reject the confident class:
\begin{equation}
a_t=\min\!\left(a,\sigma(t_2)\right),\qquad b_t=\min\!\left(b,\sigma(t_2)\right),
\label{eq:clean}
\end{equation}
where $\sigma$ is the logistic function. Scores are passed unchanged, and only
cross-class duplicates, that is, one box proposed with two class labels, are
removed. On a stream whose confident class lies above the tracker's
thresholds, the tracker therefore receives exactly its native input.

\emph{Noisy regime.} The candidates are split into bands. The primary band
$z\ge t_2$ may start and continue tracks, the extension band $t_1\le z<t_2$
may only continue tracks, and candidates below $t_1$ are discarded. Scores are
remapped through their rank $u_i$ in the empirical distribution of recent
stream scores,
\begin{equation}
\tilde{s}_i=
\begin{cases}
0.5+0.5\,u_i, & z_i\ge t_2,\\
0.1+0.4\,u_i, & t_1\le z_i<t_2,
\end{cases}
\label{eq:remap}
\end{equation}
and the tracker's association and birth thresholds are set to
$a_t=b_t=0.5$. Extension candidates can therefore reach only a tracker's
low-score stage. Duplicates are resolved with track context: a weaker
candidate that overlaps a stronger one with an intersection over union (IoU)
above one half is removed, unless a different track of frame $t-1$ claims it.
This keeps occluded people in crowds.

\subsection{Continuation rescue and motion rule}\label{sec:rescue}
A foreground candidate ($z\ge t_1$, or any emitted candidate for a
single-stage tracker when the stream has no background mode) may continue an
existing track of frame $t-1$ with IoU of at least one half while no usable candidate covers that
track and while the tracker cannot see the candidate at the operating point
passed for this frame. Such a candidate is passed with its score raised just
above the lowest threshold that the tracker uses in that frame. For a
two-stage tracker the rescue can act only on foreground candidates below the
tracker's low threshold $\ell$. For a single-stage tracker it acts on
foreground candidates below the lowest passed threshold, so it supplies the
continuation stage that the tracker lacks.

In noisy frames only, and only when an image-motion cue is available, the
matching tolerance widens with the ratio $r_t$ of the current global image
motion to its past median. The motion cue is the phase-correlation shift
\cite{reddy1996phasecorr} between consecutive grayscale frames at quarter
resolution, divided by the image diagonal. The tolerance is
\begin{equation}
m_t=\min\!\left(0.95,\;1-\frac{1-m_0}{\max(1,r_t)}\right).
\label{eq:motion}
\end{equation}

\subsection{Host contract and adapters}\label{sec:contract}
Table~\ref{tab:contracts} lists the contracts used in this paper. Each value
is a documented property of the published tracker configuration, not a fitted
quantity. An adapter for each tracker family maps the decision of the layer
onto the tracker's interface. The selected candidates and their scores replace
the detector output of the frame, and the thresholds and tolerance replace the
tracker's own for that frame. When a tracker has no overlap gate (C-TWiX) or
two detection views (TrackTrack), the adapter maps only what the tracker
exposes. In TrackTrack's second view, a detection that matches a candidate at
IoU 0.97 or more, the tracker's own rule, follows that candidate.
Post-processing that the authors apply after tracking, such as interpolation
or offline linking, is kept unchanged. The policy and all of its constants are
identical for every tracker. A tracker enters only through its published
contract values, and the adapter translates the decision without making one.

\begin{table*}[t]
\caption{Host contracts: the operating point that each published tracker
declares, as used by the layer. $\tau$ is the tracker's own published
threshold, which differs between trackers and, for TrackTrack, between
sequences.}
\label{tab:contracts}
\centering\footnotesize
\setlength{\tabcolsep}{4pt}
\begin{tabular}{@{}>{\raggedright\arraybackslash}p{30mm}lcccc>{\raggedright\arraybackslash}p{32mm}@{}}
\toprule
Tracker (setting) & Type & assoc $a$ & birth $b$ & low $\ell$ & match $m_0$ & Decision mapped to\\
\midrule
ByteTrack (official MOT17) & two-stage & 0.6 & 0.7 & 0.1 & 0.8 & thresholds, candidates, scores\\
ByteTrack (library default) & two-stage & 0.25 & 0.25 & 0.1 & 0.8 & same\\
SparseTrack & two-stage & $\tau$ & $\tau+0.1$ & 0.1 & published & same\\
BoostTrack & two-stage & $\tau$ & $\tau$ & 0.1 & $1-$IoU gate & same\\
Hybrid-SORT & two-stage & $\tau$ & $\tau$ & 0.1 & $1-$IoU gate & same\\
OC-SORT & single-stage & 0.6 & 0.6 & 0.6 & 0.7 & same\\
PD-SORT & single-stage & $\tau$ & $\tau$ & $\tau$ & $1-$IoU gate & same\\
C-TWiX (MOT17 / KITTI / DanceTrack) & single-stage & 0.5 & 0.7 / 0.5 / 0.9 & 0.5 & not mapped & candidates, scores, two score thresholds\\
TrackTrack & two-view & $\tau_{\mathrm{det}}$ & $\tau_{\mathrm{init}}$ & 0.1 & not mapped & candidates in both views, two thresholds\\
\bottomrule
\end{tabular}
\end{table*}


\subsection{Frozen constants}\label{sec:constants}
Table~\ref{tab:constants} lists every constant of the layer. The same values
are used for every tracker, detector, and dataset. The band thresholds $t_1$
and $t_2$ are not constants: they are recomputed online for every frame. The
duplicate and rescue overlap of one half follows the matching rule of MOT
evaluation, in which two boxes with IoU above 0.5 cannot both match distinct
objects. The remap ranges of Eq.~\eqref{eq:remap} place the two bands on the
common association and birth threshold of 0.5 and the low-score threshold of
0.1. The cap of 0.95 keeps a minimum overlap of 0.05 for every match. The
window, the motion history, and the remap memory are declared defaults that
were kept fixed in all development stages.

\begin{table}[t]
\caption{Frozen constants of \scl{}.}
\label{tab:constants}
\centering\footnotesize
\begin{tabular}{@{}>{\raggedright\arraybackslash}p{39mm}>{\raggedright\arraybackslash}p{31mm}@{}}
\toprule
Constant & Value\\
\midrule
Statistics window $W$ & 10 frames\\
Regime rule & clean if median $\rho\ge0.5$ over non-cold frames so far\\
Background check & $\rho=1$ if fewer candidates below $t_1$ than above\\
Noisy score ranges & primary $[0.5,1)$, extension $[0.1,0.5)$\\
Noisy association and birth thresholds & 0.5\\
Remap memory & 20 frames, one every 10 frames\\
Duplicate and rescue IoU & 0.5\\
Rescued score & lowest used threshold $+\,0.001$\\
Motion cue & phase correlation, quarter resolution\\
Motion history, warm-up & 100, 5 frames\\
Tolerance cap & 0.95\\
\bottomrule
\end{tabular}
\end{table}

\subsection{Per-frame procedure and cost}\label{sec:cost}
Algorithm~\ref{alg:layer} gives the procedure for one frame. Per frame the
layer sorts the pooled logits of at most $W$ frames for the two Otsu splits,
computes candidate overlaps with the other candidates and with the tracks of
the previous frame, and, when the motion cue is used, computes one phase
correlation on a quarter-resolution image. With $n$ candidates per frame and
$k$ tracks, the cost is $O(Wn\log(Wn)+n^2+nk)$ plus the fixed image cost. The
state is a ring buffer of $W$ frames of logits, the per-frame $\rho$ values,
the remap memory, and a history of 100 motion values.

\begin{algorithm}[t]
\caption{\scl{} for frame $t$}\label{alg:layer}
\begin{algorithmic}[1]
\Require candidates $D_t$, tracks $\mathcal{T}_{t-1}$, cue $g_t$, contract $(a,b,\ell,m_0)$
\State $z_i\gets$ logit of $s_i$ for every candidate
\If{no valid split exists}
  \State regime $\gets$ cold; $(a_t,b_t)\gets(a,b)$
\Else
  \State $(t_1,t_2)\gets$ nested Otsu on the window
  \State $\rho_t\gets$ Eq.~\eqref{eq:rho} with the background check
  \State regime $\gets$ clean if $\mathrm{median}(\rho)\ge0.5$, else noisy
\EndIf
\If{regime is cold or clean}
  \State remove cross-class duplicates; keep raw scores
  \If{clean} set $(a_t,b_t)$ by Eq.~\eqref{eq:clean}
  \EndIf
\Else
  \State remove duplicates with track context
  \State drop $z<t_1$; remap scores by Eq.~\eqref{eq:remap}
  \State $a_t\gets0.5$; $b_t\gets0.5$
\EndIf
\State rescue track continuations (not in cold frames)
\State $m_t\gets$ Eq.~\eqref{eq:motion} if noisy, else $m_0$
\State pass $K_t$, $\tilde{s}$, $a_t$, $b_t$, $m_t$ to the tracker
\State add the logits of frame $t$ to the window; add $g_t$ to the history
\State store the tracks of frame $t$ after the tracker update
\end{algorithmic}
\end{algorithm}
